\heading{数学物理方法（下）-21春-期末}
\noteinfo{题目来源于赛艇，答案由AI生成。}

\begin{question}[积分题]

计算
\[
\int_0^x \cos(x-t)J_0(t)\,dt.
\]

\begin{solution}

这是卷积。记
\[
F(x)=\int_0^x \cos(x-t)J_0(t)\,dt.
\]
作 Laplace 变换：
\[
\mathcal L\{\cos x\}=\frac{s}{s^2+1},
\qquad
\mathcal L\{J_0(x)\}=\frac{1}{\sqrt{s^2+1}}.
\]
于是
\[
\mathcal L\{F\}=\frac{s}{(s^2+1)^{3/2}}.
\]
又由 Bessel 函数恒等式
\[
\frac{d}{dx}[xJ_1(x)]=xJ_0(x)
\]
可知
\[
\mathcal L\{xJ_0(x)\}=\frac{s}{(s^2+1)^{3/2}}.
\]
所以
\ansmath{\int_0^x \cos(x-t)J_0(t)\,dt=xJ_0(x).}

\end{solution}
\end{question}

\begin{question}[圆盘内热方程]

求解定解问题
\[
\begin{cases}
\dfrac{\partial u}{\partial t}-\kappa\nabla^2u=y,\\
u|_{x^2+y^2=a^2}=x^2-y^2,\quad t>0,\\
u|_{t=0}=0,\quad x^2+y^2<a^2.
\end{cases}
\]

\begin{solution}

令 $u=w+v$，其中 $w$ 是与时间无关的稳态解，满足
\[
-\kappa \nabla^2 w=y,
\qquad w|_{r=a}=x^2-y^2.
\]
先取边界项的调和延拓
\[
 w_1=x^2-y^2=r^2\cos 2\theta.
\]
再取零边界特解
\[
 w_2=-\frac{(r^2-a^2)y}{8\kappa}=-\frac{(r^2-a^2)r\sin\theta}{8\kappa},
\]
因为在二维中有
\[
\nabla^2(r^2 y)=8y.
\]
故
\[
 w(x,y)=x^2-y^2-\frac{(x^2+y^2-a^2)y}{8\kappa}.
\]
于是 $v=u-w$ 满足齐次热方程、零边界和初值 $v(x,y,0)=-w(x,y)$。它可按圆盘上的 Bessel 本征函数展开：
\[
 v(r,\theta,t)=
 \sum_{n\ge1}A_n e^{-\kappa j_{1,n}^2 t/a^2}J_1\!\Bigl(\frac{j_{1,n}r}{a}\Bigr)\sin\theta
 +
 \sum_{n\ge1}B_n e^{-\kappa j_{2,n}^2 t/a^2}J_2\!\Bigl(\frac{j_{2,n}r}{a}\Bigr)\cos2\theta,
\]
系数由 $-w$ 的正交展开确定。因此标准答案可写为
\ansmath{u=w+v,\quad w=x^2-y^2-\dfrac{(x^2+y^2-a^2)y}{8\kappa}.}

\end{solution}
\end{question}

\begin{question}[本征值问题的自伴性]

回忆版题面中第三题写成
\[
\begin{cases}
y''+\lambda y=0,\qquad 0<x<l,\\
y(l)=\alpha y(0),\qquad y'(l)=\beta y(l).
\end{cases}
\]
并询问何时自伴。

\begin{solution}

若边界条件确为
\[
 y(l)=\alpha y(0),\qquad y'(l)=\beta y(l),
\]
则 Green 边界型
\[
 [\bar y z'-\bar y' z]_{0}^{l}
\]
在代入边界条件后，仍会保留端点 $x=0$ 处的 $y'(0),z'(0)$ 自由项，无法对所有满足边界条件的 $y,z$ 恒为零。因此这组边界条件在非退化情形下\emph{不是}自伴边界条件。

也就是说，按当前补正后的题面，自伴性的答案应为：\textbf{一般不存在使其成为自伴边界条件的非退化参数取值。}

\ansmath{[\bar y z'-\bar y' z]_{0}^{l}}

\end{solution}
\end{question}

\begin{question}[均匀带电圆盘的电荷密度]

有一总电量为 $Q$、半径为 $a$ 的均匀带电圆盘，分别给出直角坐标、柱坐标、球坐标下空间的电荷密度分布函数。

\begin{solution}

圆盘面电荷密度为
\[
\sigma=\frac{Q}{\pi a^2}.
\]
故在三种坐标下可写为：
\begin{enumerate}[label=(\arabic*),leftmargin=2.6em,itemsep=0.25em,topsep=0.2em]
\item 直角坐标：
\[
\rho(x,y,z)=\frac{Q}{\pi a^2}\,\delta(z)\,H(a^2-x^2-y^2).
\]
\item 柱坐标：
\[
\rho(r,\varphi,z)=\frac{Q}{\pi a^2}\,\delta(z)\,H(a-r).
\]
\item 球坐标：由于 $z=r\cos\theta$，且平面 $z=0$ 对应 $\theta=\pi/2$，因此
\[
\rho(r,\theta,\varphi)=\frac{Q}{\pi a^2}\,\frac{\delta(\theta-\pi/2)}{r}\,H(a-r).
\]
\end{enumerate}
其中 $H$ 为 Heaviside 阶跃函数。

\ansmath{\rho(r,\varphi,z)=\frac{Q}{\pi a^2}\,\delta(z)\,H(a-r)\quad \rho(r,\theta,\varphi)=\frac{Q}{\pi a^2}\,\frac{\delta(\theta-\pi/2)}{r}\,H(a-r)}

\end{solution}
\end{question}

\begin{question}[三维无界空间波动方程的 Green 函数]

求解三维无界空间中
\[
\begin{cases}
\dfrac{\partial^2u}{\partial t^2}-c^2\nabla^2u=f(\mathbf r),\\
u|_{t=0}=g(\mathbf r),\qquad \left.\dfrac{\partial u}{\partial t}\right|_{t=0}=h(\mathbf r)
\end{cases}
\]
对应的 Green 函数。

\begin{solution}

三维波动算子的迟滞 Green 函数为
\ansmath{G(\mathbf r,t;\mathbf r',t')=
\frac{\delta\!\left(t-t'-\dfrac{|\mathbf r-\mathbf r'|}{c}\right)}{4\pi c^2|\mathbf r-\mathbf r'|}
\,H(t-t').}
因此 Kirchhoff 公式写成
\[
 u(\mathbf r,t)=
 \iiint G(\mathbf r,t;\mathbf r',0)h(\mathbf r')\,d^3\mathbf r'
 +\frac{\partial}{\partial t}\iiint G(\mathbf r,t;\mathbf r',0)g(\mathbf r')\,d^3\mathbf r'
 +\iiint\!\int_0^t G(\mathbf r,t;\mathbf r',\tau)f(\mathbf r',\tau)\,d\tau\,d^3\mathbf r'.
\]
若 $f$ 与时间无关，上式第三项再对 $\tau$ 直接积分即可。

\end{solution}
\end{question}

\clearpage

